\capitulo{1}{Introducción}
El presente Trabajo Fin de Grado consiste en el desarrollo de un sistema SIEM (Security Information and Event Management) ligero orientado a la monitorización de eventos de seguridad dentro de un entorno simulado de alta seguridad basado en un polvorín militar.

Actualmente, la protección de infraestructuras críticas tiene una gran importancia desde múltiples perspectivas de seguridad, incluyendo la seguridad física, la seguridad lógica, la ciberseguridad y la protección de los sistemas de información que soportan la operación de estas instalaciones.

En este tipo de instalaciones resulta necesario disponer de mecanismos que permitan supervisar de forma continua los diferentes elementos de seguridad y detectar posibles incidencias o comportamientos anómalos.

Con el fin de simular una situación lo más cercana posible a un entorno real, se ha diseñado un polvorín militar compuesto por diferentes depósitos y sistemas de protección. Entre ellos se incluyen controles de acceso mediante tarjeta, sensores magnéticos de apertura de puertas, sensores de movimiento, sensores ambientales de temperatura y humedad, cámaras de vigilancia y sistemas de respaldo eléctrico.

La elección de un polvorín militar permite disponer de un escenario donde la supervisión continua y la detección temprana de incidencias resultan especialmente relevantes debido a la naturaleza de los materiales almacenados.

A partir de la información generada por estos elementos, el sistema desarrollado es capaz de recoger eventos, almacenarlos en una base de datos, aplicar reglas de correlación y mostrar la información mediante una interfaz web de monitorización. 

Además de la monitorización de eventos, el sistema incorpora funcionalidades de supervisión de sensores, detección de anomalías, generación de alertas y visualización de información mediante una interfaz web. Todo ello se ha desarrollado utilizando tecnologías de código abierto y buscando una solución sencilla, modular y fácilmente ampliable.

Aunque el entorno utilizado durante el desarrollo es simulado, la arquitectura propuesta puede servir como base para futuros desarrollos orientados a la protección de infraestructuras críticas o entornos con requisitos especiales de seguridad.

La protección de instalaciones que albergan material sensible requiere la aplicación de medidas de seguridad física y lógica que permitan detectar accesos no autorizados, alteraciones ambientales y posibles incidencias operativas. En este contexto, organismos nacionales e internacionales han desarrollado marcos normativos y recomendaciones orientadas a la protección de infraestructuras críticas y activos estratégicos.

El presente trabajo toma como referencia la Ley 8/2011, de 28 de abril \cite{ley811}, por la que se establecen medidas para la protección de las infraestructuras críticas, así como recomendaciones de seguridad empleadas en entornos de defensa recogidas en documentación de la Organización del Tratado del Atlántico Norte (OTAN) \cite{stanag4440} \cite{aastp5} \cite{aastp1} . Asimismo, se consideran buenas prácticas derivadas de estándares internacionales de seguridad de la información, como la norma ISO/IEC 27001 \cite{iso27001}.
 

